\input{sol_pre.tex}
\input{q_2026-03-07_2.tex}
\begin{document}
\input{sol_head.tex}

%================= 1
\Savol{1}
$24^3=2^9\cdot3^3$, \ $6^{12}=2^{12}\cdot3^{12}$, \ $20^3=2^6\cdot5^3$ $\Rightarrow$ $m=2^{27}\cdot3^{15}\cdot5^3$, \ $m^2=2^{54}\cdot3^{30}\cdot5^6$.\\
$\dfrac{m^2}{2^n}$ natural son bölişi uçun $n\le54$.
\Javob{A) 54}

%================= 2
\Savol{2}
$x\in\left[\frac12;2\right]$ da: $3x-1>0$, $x>0$, $x-2\le0$. \ $|3x-1|-|x|=2x-1\ge0$:
\[(2x-1)-(2-x)=3x-3=3(x-1).\]
\Javob{B) $3(x-1)$}

%================= 3
\Savol{3}
Katta kenguru 2 marta sakraganda $2\cdot2=4$ m, kiçigi şu vaqtda $3\cdot1=3$ m yuradi: har «sikl»da masofa 1 m ga qisqaradi.\\
30 m qisqarişi uçun 30 sikl kerak. Katta kenguru: $30\cdot4=120$ m.
\Javob{A) 120}

%================= 4
\Savol{4}
$7\frac12:\frac74=\frac{15}2:\frac74=30:7$. \ Birinçi $30k$, ikkinçi $7k$, uçinçi $43\frac13\%\cdot30k=\frac{13}{30}\cdot30k=13k$.
\[30k+7k+13k=50k=4\,080\,000\ \Rightarrow\ k=81\,600,\qquad 13k=1\,060\,800.\]
\Javob{A) 1\,060\,800}

%================= 5
\Savol{5}
Daraja körsatkiçlarini soliştiramiz: \ A) $2^{2^2}=16$; \ B) $2^{22}=4\,194\,304$; \ C) $222$; \ D) $22^2=484$.\\
Asos bir xil ($2>1$), şuning uçun eng kattasi körsatkiçi eng katta bölgan $2^{2^{22}}$.
\Javob{B) $2^{2^{22}}$}

%================= 6
\Savol{6}
$\sqrt2\cdot\sqrt{3-\sqrt5}=\sqrt{6-2\sqrt5}=\sqrt5-1$:
\[(\sqrt5-1)^2(3+\sqrt5)=(6-2\sqrt5)(3+\sqrt5)=18+6\sqrt5-6\sqrt5-10=8.\]
\Javob{D) 8}

%================= 7
\Savol{7}
\[\sqrt[3]{\frac{125x^{\frac65}}{54}}=\frac{5x^{\frac25}}{\sqrt[3]{27\cdot2}}=\frac{5x^{\frac25}}{3\sqrt[3]2}=\frac{5\sqrt[3]4\,x^{\frac25}}{3\cdot2}=\frac{5\sqrt[3]4\,x^{\frac25}}6.\]
\Javob{A) $\frac{5\sqrt[3]4\,x^{\frac25}}6$}

%================= 8
\Savol{8}
$a_7=5+2d$, \ $a_{10}=5+5d$:
\[a_7a_{10}=10d^2+35d+25,\qquad d_{\min}=-\frac{35}{2\cdot10}=-1{,}75.\]
\Javob{D) $-1{,}75$}

%================= 9
\Savol{9}
$b_2^2=b_1b_3$: \ $a_4^2=a_2a_6=(a_4-2d)(a_4+2d)=a_4^2-4d^2\Rightarrow d=0$.\\
Demak barça hadlar teng: $3a_1=16\Rightarrow a_1=\frac{16}3$. \quad $n\cdot\dfrac{16}3=176\Rightarrow n=33$.
\Javob{C) 33}

%================= 10
\Savol{10}
$(2a+b)^2-3a(a+b)=a^2+ab+b^2$, \ $a^3-b^3=(a-b)(a^2+ab+b^2)$ $\Rightarrow$ birinçi kasr $a-b$.\\
$(a^2+b)(b^2+a)-ab(ab+1)=a^2b^2+a^3+b^3+ab-a^2b^2-ab=a^3+b^3$ $\Rightarrow$ ikkinçi kasr $a+b$.
\[(a-b)+(a+b)=2a.\]
\Javob{D) $2a$}

%================= 11
\Savol{11}
$a^3=1-a$ $\Rightarrow$ $a^4=a-a^2$, \ $a^5=a^2\cdot a^3=a^2-a^3=a^2+a-1$.
\[\frac{(a-a^2)+(1-a)+a^2+9}{(a^2+a-1)-a^2-a+6}=\frac{10}5=2.\]
\Javob{A) 2}

%================= 12
\Savol{12}
\[3\sin10^\circ-\sqrt3\cos10^\circ=2\sqrt3\left(\tfrac{\sqrt3}2\sin10^\circ-\tfrac12\cos10^\circ\right)=2\sqrt3\sin(10^\circ-30^\circ)=-2\sqrt3\sin20^\circ.\]
$3\cos70^\circ=3\sin20^\circ$: \quad $\dfrac{-2\sqrt3\sin20^\circ}{3\sin20^\circ}=-\dfrac{2\sqrt3}3=-\dfrac2{\sqrt3}$.
\Javob{B) $-\frac2{\sqrt3}$}

%================= 13
\Savol{13}
$\alpha+\beta+\gamma=\pi$ da: \ $\sin2\alpha+\sin2\beta=2\sin(\alpha+\beta)\cos(\alpha-\beta)=2\sin\gamma\cos(\alpha-\beta)$, \ $\sin2\gamma=-2\sin\gamma\cos(\alpha+\beta)$.
\[\sin2\alpha+\sin2\beta+\sin2\gamma=2\sin\gamma\bigl(\cos(\alpha-\beta)-\cos(\alpha+\beta)\bigr)=4\sin\alpha\sin\beta\sin\gamma=4\cdot\frac12=2.\]
\Javob{A) 2}

%================= 14
\Savol{14}
$77^x-121^x=11^x(7^x-11^x)>0\Rightarrow 7^x>11^x\Rightarrow x<0$; unda surat musbat. Kvadratga kötaramiz:
\[(7^x-11^x)^2=11^x(7^x-11^x)\ \Rightarrow\ 7^x-11^x=11^x\ \Rightarrow\ \left(\frac7{11}\right)^x=2\ \Rightarrow\ x=\log_{\frac7{11}}2<0.\]
Yagona ildiz.
\Javob{B) 1}

%================= 15
\Savol{15}
Öng tomon: $\log_{4-x}1=0$, mavjud bölişi uçun $4-x>0$, $4-x\ne1$, ya'ni $x<4$, $x\ne3$.
\[\log_{\frac12}\left(2x-5{,}5\right)\ge0\ \Rightarrow\ 0<2x-5{,}5\le1\ \Rightarrow\ 2{,}75<x\le3{,}25.\]
Bu oraliqdagi yagona butun son $x=3$, lekin $x\ne3$. Butun yeçim yöq.
\Javob{D) Butun yeçimga ega emas}

%================= 16
\Savol{16}
$D=9-4a^2\ge0\Rightarrow |a|\le1{,}5$. \ Butun qiymatlar: $-1;\ 0;\ 1$ — 3 ta.
\Javob{C) 3}

%================= 17
\Savol{17}
$\sqrt[3]{9-x^2}=-\sqrt[3]{(x+3)(x-3)}$. \ $u=\sqrt[3]{x+3}$, \ $v=\sqrt[3]{x-3}$:
\[u^2-6v^2=-uv\ \Rightarrow\ u^2+uv-6v^2=0\ \Rightarrow\ (u+3v)(u-2v)=0.\]
$u=-3v$: $x+3=-27(x-3)\Rightarrow x=\frac{78}{28}=\frac{39}{14}$. \quad $u=2v$: $x+3=8(x-3)\Rightarrow x=\frac{27}7=\frac{54}{14}$.
\[\frac{39}{14}+\frac{54}{14}=\frac{93}{14}=6\frac9{14}.\]
\Javob{C) $6\frac9{14}$}

%================= 18
\Savol{18}
$x^2-7x+10=(x-2)(x-5)$, \ $\dfrac8{5-x}=-\dfrac{8(x-2)}{(x-2)(x-5)}$:
\[\frac{x^2+3x+46+8x-16}{(x-2)(x-5)}\le0\ \Rightarrow\ \frac{(x+5)(x+6)}{(x-2)(x-5)}\le0\ \Rightarrow\ x\in[-6;-5]\cup(2;5).\]
Butun yeçimlar: $-6,-5,3,4$; \ yiğindisi $-4$.
\Javob{D) $-4$}

%================= 19
\Savol{19}
Aniqlaniş sohasi: $x\ge11$ va $x\le12$, ya'ni $[11;12]$. Ikkala ildiz bir vaqtda nolga teng bölmaydi ($x=11$ va $x=12$ bir vaqtda emas), şuning uçun yiğindi har doim musbat.
\Javob{B) $[11;12]$}

%================= 20
\Savol{20}
$x+3\ge0$, \ $6-x\ge0$, \ $(x-5)^2\ne0$ $\Rightarrow$ $x\in[-3;6]$, $x\ne5$.
\[(-3)+(-2)+\dots+6-5=15-5=10.\]
\Javob{A) 10}

%================= 21
\Savol{21}
$x^3+1=(x+1)(x^2-x+1)$. \ Sodda kasrlarga yoyamiz:
\[\frac1{x^3+1}=\frac13\cdot\frac1{x+1}-\frac13\cdot\frac{x-2}{x^2-x+1},\qquad \frac{x-2}{x^2-x+1}=\frac12\cdot\frac{2x-1}{x^2-x+1}-\frac32\cdot\frac1{\left(x-\frac12\right)^2+\frac34}.\]
\[\int\frac{dx}{x^3+1}=\frac13\ln|x+1|-\frac16\ln(x^2-x+1)+\frac12\cdot\frac2{\sqrt3}\arctg\frac{2x-1}{\sqrt3}+C.\]
\Javob{C}

%================= 22
\Savol{22}
$S_n(x)=x+x^2+\dots+x^n$: \ $S_n(1)=n$, \ $S_n'(1)=1+2+\dots+n=\dfrac{n(n+1)}2$.
\[f'(1)=\frac{S_{2024}'(1)S_{1012}(1)-S_{2024}(1)S_{1012}'(1)}{S_{1012}(1)^2}=\frac{\frac{2024\cdot1012}2\,(2025-1013)}{1012^2}=\frac{2024}2=1012.\]
\Javob{D) 1012}

%================= 23
\Savol{23}
\fig{\begin{tikzpicture}[scale=1.1]
\def\R{2.414}
\draw[thick] (0,0)--(\R,0) arc (0:90:\R)--cycle;
\draw[thick] (1,1) circle (1);
\draw[axs] (0,0)--(45:\R); \draw[ax] (1,1)--(1,0);
\fill (1,1) circle (1.5pt) node[left]{\scriptsize$O_1$};
\node[below left] at (0,0){\small$O$}; \node[right] at (\R,0){\small$A$}; \node[left] at (0,\R){\small$B$};
\node[aux,right] at (1,0.5){\scriptsize$r$}; \node[aux] at (1.25,1.55){\scriptsize$r$}; \node[aux] at (0.5,0.75){\scriptsize$r\sqrt2$};
\end{tikzpicture}}
Aylana $OA$ va $OB$ ga urinadi, demak $O_1$ töğri burçak bissektrisasida, $OO_1=r\sqrt2$.\\
Çorak aylanaga uriniş nuqtasi $OO_1$ davomida: \ $r\sqrt2+r=1+\sqrt2\Rightarrow r=1$.
\Javob{B) 1}

%================= 24
\Savol{24}
$x^2-ax+3=2x-1\Rightarrow x^2-(a+2)x+4=0$. Bitta umumiy nuqta: $D=(a+2)^2-16=0$.
\[a+2=\pm4\ \Rightarrow\ a=2\ \text{yoki}\ a=-6;\qquad 2+(-6)=-4.\]
\Javob{B) $-4$}

%================= 25
\Savol{25}
Uçinçi burçak $180^\circ-64^\circ-76^\circ=40^\circ$. \ $\alpha$ va $\beta$ burçaklar bissektrisalari orasidagi ötkir burçak $\dfrac{\alpha+\beta}2$ (uçburçakdan $180^\circ-\frac\alpha2-\frac\beta2$ ning qöşnisi):
\[\frac{64^\circ+76^\circ}2=70^\circ,\quad \frac{64^\circ+40^\circ}2=52^\circ,\quad \frac{76^\circ+40^\circ}2=58^\circ.\]
\Javob{A) $52^\circ$}

%================= 26
\Savol{26}
$a=BC=13$, $b=AC=15$, $c=AB=8$, \ $p=18$, \ $S=\sqrt{18\cdot5\cdot3\cdot10}=30\sqrt3$, \ $r=\dfrac Sp=\dfrac{5\sqrt3}3$.\\
Urinmalar: $A$ dan $p-a=5$, $B$ dan $p-b=3$, $C$ dan $p-c=10$. \ $LA^2=r^2+5^2=\frac{100}3$, \ $LB^2=r^2+9=\frac{52}3$, \ $LC^2=r^2+100=\frac{325}3$.\\
Leybnits teoremasi ($M$ — medianalar kesişiş nuqtasi): \ $LA^2+LB^2+LC^2=3LM^2+\dfrac{a^2+b^2+c^2}3$:
\[159=3LM^2+\frac{458}3\ \Rightarrow\ LM^2=\frac{19}9\ \Rightarrow\ LM=\frac{\sqrt{19}}3.\]
\Javob{B) $\frac{\sqrt{19}}3$}

%================= 27
\Savol{27}
Töğri burçaklar tufayli $5$, $4$, $3$ uzunlikdagi kesmalar özaro parallel, $2$ va $2$ kesmalar ularga perpendikulyar (va bir tomonga yönalgan).\\
Şuning uçun $B$ dan $D$ ga köçiş: bir yönalişda $5+4+3=12$, unga perpendikulyar yönalişda $2+2=4$:
\[BD^2=12^2+4^2=160,\qquad S_{\text{kv}}=\frac{BD^2}2=80.\]
\Javob{B) 80}

%================= 28
\Savol{28}
Aylana radiusi $r$. Içki çizilgan oltiburçak tomoni $a_1=r$; taşqi çizilganining apofemasi $r$, tomoni $a_2=\dfrac{2r}{\sqrt3}$.\\
Öxşaş şakllar yuzlari nisbati — tomonlar nisbatining kvadrati: \ $\left(\dfrac{a_2}{a_1}\right)^2=\dfrac43$.
\Javob{B) $\frac43$}

%================= 29
\Savol{29}
$B(x;y;0)$, \ $\overrightarrow{AB}=(x-6;\ y-4;\ 2)\parallel(4;-3;1)$ $\Rightarrow$ $k=2$: \ $\overrightarrow{AB}=(8;-6;2)$.\\
Koordinatalar yiğindisi: $8-6+2=4$.
\Javob{A) 4}

%================= 30
\Savol{30}
Öq kesimi — tomonlari $2r$ va $h$ töğri törtburçak: \ $(4r)^2=(2r)^2+h^2\Rightarrow h=2\sqrt3\,r$.\\
$2\pi rh=4\sqrt3\pi r^2=12\sqrt3\pi\Rightarrow r^2=3$, \ $h=6$. \quad $V=\pi r^2h=18\pi$.
\Javob{C) $18\pi$}

%================= 31
\Savol{31}
Böyalgan soha $A$ va $C$ ning umumiy qismi, lekin $B$ ga kirmaydi: $(A\cap C)\setminus B$.
\Javob{A) $(A\cap C)\setminus B$}

%================= 32
\Savol{32}
$\overline{abc}$: $a\ge1$. \ $a'=a-1\ge0$: \ $a'+b+c=8$ ning manfiymas butun yeçimlari soni $C_{10}^2=45$ (barça raqamlar $\le8$, çeklov buzilmaydi).
\Javob{B) 45}

%================= 33-35
\newpage
\Savol{33}
\fig{\begin{tikzpicture}[scale=0.8]
\coordinate (O) at (0,0); \coordinate (A) at (-2,-0.5); \coordinate (B) at (0.6,-1); \coordinate (S) at (0,3.2); \coordinate (K) at (-0.7,-0.75);
\draw[thick] (0,0) ellipse (2.4 and 0.8);
\draw[thick] (-2.4,0)--(S)--(2.4,0);
\draw[thick,fill=blue!20,fill opacity=0.6] (A)--(B)--(S)--cycle;
\draw[ax] (S)--(O); \draw[axs] (O)--(B); \draw[ax] (S)--(K)--(O);
\fill (O) circle (1.4pt) node[right]{\scriptsize$O$};
\node[left] at (A){\small$A$}; \node[below] at (B){\small$B$}; \node[above] at (S){\small$S$}; \node[below] at (K){\scriptsize$K$};
\node[aux,right] at (0.3,-0.4){\scriptsize$R$};
\end{tikzpicture}}
Kvadrat tomoni 4, uning diagonali asos diametri: $2R=4\sqrt2\Rightarrow R=2\sqrt2$.\\
$\triangle ASB$: $SA=SB$ (yasovçilar), $\angle ASB=60^\circ$ $\Rightarrow$ teng tomonli, $SA=AB=4$.

\textbf{33.} \ $h=SO=\sqrt{SA^2-R^2}=\sqrt{16-8}=2\sqrt2$.

\textbf{34.} \ $S_{ABS}=\dfrac{\sqrt3}4\cdot4^2=4\sqrt3$.

\textbf{35.} \ $V=\dfrac13\pi R^2h=\dfrac13\cdot3\cdot8\cdot2\sqrt2=16\sqrt2$.
\Javob{33 – C; \ 34 – E; \ 35 – A}

%================= 36
\newpage
\Savol{36}
$u=x-y$, \ $xy=2u$. \ $x^2+y^2=u^2+2xy=u^2+4u$:
\[x^4+y^4=(x^2+y^2)^2-2(xy)^2=(u^2+4u)^2-8u^2=17u.\]
$u=0$: $xy=0$ va $x=y$ $\Rightarrow$ $(0;0)$. \\
$u\ne0$: $u(u+4)^2-8u=17\Rightarrow u^3+8u^2+8u-17=0\Rightarrow (u-1)(u^2+9u+17)=0$.\\
$x$ va $-y$ — $t^2-ut-2u=0$ ning ildizlari (yiğindisi $u$, köpaytmasi $-2u$); haqiqiy bölişi uçun $u^2+8u\ge0$, ya'ni $u\ge0$ yoki $u\le-8$.\\
$u^2+9u+17=0$ ildizlari $\frac{-9\pm\sqrt{13}}2\in(-8;0)$ — mos emas.\\
$u=1$: $t^2-t-2=0\Rightarrow t=2$ yoki $t=-1$ $\Rightarrow$ $(2;1)$, \ $(-1;-2)$.

a) Yeçimlar: $(0;0)$, $(2;1)$, $(-1;-2)$ — 3 ta.

b) $x+y$: $0;\ 3;\ -3$. Eng kiçigi $-3$.
\Javob{a) 3; \quad b) $-3$}

%================= 37
\newpage
\Savol{37}
$\sin^3x=\dfrac{3\sin x-\sin3x}4$, \ $\cos^3x=\dfrac{3\cos x+\cos3x}4$:
\[\frac{\cos3x(3\sin x-\sin3x)+\sin3x(3\cos x+\cos3x)}4=\frac34\left(\sin x\cos3x+\cos x\sin3x\right)=\frac34\sin4x.\]
\[\frac34\sin4x=\frac34\ \Rightarrow\ \sin4x=1\ \Rightarrow\ x=\frac\pi8+\frac{\pi k}2,\ k\in\mathbb Z.\]
a) Eng kiçik musbat ildiz ($k=0$): $\dfrac\pi8$.

b) $-2\pi\le\dfrac\pi8+\dfrac{\pi k}2\le2\pi\Rightarrow -4{,}25\le k\le3{,}75\Rightarrow k=-4,\dots,3$ — 8 ta.
\Javob{a) $\frac\pi8$; \quad b) 8}

%================= 38
\newpage
\Savol{38}
a) $\cos\sqrt2x$ ning davri $T_1=\dfrac{2\pi}{\sqrt2}=\sqrt2\pi$, \ $\cos\dfrac x{\sqrt2}$ niki $T_2=2\sqrt2\pi=2T_1$.\\
Umumiy eng kiçik davr $T=2\sqrt2\pi$ (ikkinçi qöşiluvçining davri kiçikroq emas, $T_1$ esa unga davr emas).

b) $t=\dfrac x{\sqrt2}$: \ $\sqrt2x=2t$,
\[f=\cos2t+\cos t=2\cos^2t+\cos t-1=2\left(\cos t+\frac14\right)^2-\frac98.\]
$\cos t=-\dfrac14$ da eng kiçik qiymat $-\dfrac98$.
\Javob{a) $2\sqrt2\pi$; \quad b) $-\frac98$}

%================= 39
\newpage
\Savol{39}
Grafikdan: $(-1;5)$ da $f'(x)<0$ (tölqinlar $Ox$ ga yetmaydi), $f'$ işorasi $x=5$ da «$-$» dan «$+$» ga, $x=9$ da «$+$» dan «$-$» ga, $x=11$ da «$-$» dan «$+$» ga özgaradi.

a) Ekstremum nuqtalari: $x=5$ (min), $x=9$ (max), $x=11$ (min) — 3 ta.

b) $[5;11]$ da $f$ $[5;9]$ da ösadi, $[9;11]$ da kamayadi $\Rightarrow$ eng katta qiymat $x=9$ da.
\Javob{a) 3; \quad b) 9}

%================= 40
\newpage
\Savol{40}
a) $f(x)=g(x)$: $x^2-5x+4=0\Rightarrow x=1;\ 4$.
\[S_1=\int_1^4\bigl((x+1)-(x^2-4x+5)\bigr)dx=\int_1^4(-x^2+5x-4)dx=\left(-\frac{x^3}3+\frac{5x^2}2-4x\right)\Big|_1^4=\frac83+\frac{11}6=\frac92.\]
b) $f(x)=h(x)$: $x^2-5x=0\Rightarrow x=0;\ 5$.
\[\int_0^5\bigl((x+5)-(x^2-4x+5)\bigr)dx=\int_0^5(5x-x^2)dx=\frac{125}2-\frac{125}3=\frac{125}6.\]
$S_2$ — $h$ va $f$ orasidagi sohadan $S_1$ ni olib taşlagani: \ $S_2=\dfrac{125}6-\dfrac92=\dfrac{98}6=\dfrac{49}3$.
\Javob{a) $\frac92$; \quad b) $\frac{49}3$}

%================= 41
\newpage
\Savol{41}
$\beta=\angle APB$, \ $\angle APC=360^\circ-120^\circ-\beta=240^\circ-\beta$.\\
$\triangle ABP$: $\angle ABP=165^\circ-\beta$; \quad $\triangle ACP$: $\angle ACP=180^\circ-15^\circ-(240^\circ-\beta)=\beta-75^\circ$.\\
Sinuslar teoremasi: \ $\dfrac{AP}{\sin\angle ABP}=\dfrac{BP}{\sin15^\circ}$, \ $\dfrac{AP}{\sin\angle ACP}=\dfrac{PC}{\sin15^\circ}$ \ $\Rightarrow$ \ $\sqrt3\sin(165^\circ-\beta)=\sin(\beta-75^\circ)$.

a) $\varphi=\beta-75^\circ$ desak, $165^\circ-\beta=90^\circ-\varphi$:
\[\sqrt3\cos\varphi=\sin\varphi\ \Rightarrow\ \tg\varphi=\sqrt3\ \Rightarrow\ \varphi=60^\circ,\ \ \beta=135^\circ,\ \ \sin\beta=\frac{\sqrt2}2.\]
b) $\triangle APC$: $\angle ACP=60^\circ$, $\angle APC=105^\circ$, $\angle PAC=15^\circ$. Sinuslar teoremasi:
\[AC=\frac{PC\cdot\sin105^\circ}{\sin15^\circ}=\frac{\cos15^\circ}{\sin15^\circ}=\ctg15^\circ=2+\sqrt3.\]
\Javob{a) $\frac{\sqrt2}2$; \quad b) $2+\sqrt3$}

%================= 42
\newpage
\Savol{42}
$AB=x$, $AD=BC=3x$, \ $\alpha=\angle EAD$, \ $\angle BAE=90^\circ-\alpha$. Kosinuslar teoremasi:
\[\triangle ABE:\ 2=x^2+1-2x\sin\alpha\ \Rightarrow\ \sin\alpha=\frac{x^2-1}{2x};\qquad \triangle AED:\ 16=9x^2+1-6x\cos\alpha\ \Rightarrow\ \cos\alpha=\frac{9x^2-15}{6x}.\]
$\sin^2\alpha+\cos^2\alpha=1$, \ $s=x^2$: \ $9(s-1)^2+(9s-15)^2=36s\Rightarrow 90s^2-324s+234=0\Rightarrow 5s^2-18s+13=0$.\\
$s=1$ da $\sin\alpha=0$ — $E$ içkarida emas. Demak $s=x^2=\dfrac{13}5$.

a) $\cos\alpha=\dfrac{9\cdot\frac{13}5-15}{6x}=\dfrac{42/5}{6x}=\dfrac7{5x}=\dfrac7{5\sqrt{13/5}}=\dfrac7{\sqrt{65}}$.

b) $S_{ABCD}=x\cdot3x=3\cdot\dfrac{13}5=\dfrac{39}5$.
\Javob{a) $\frac7{\sqrt{65}}$; \quad b) $\frac{39}5$}

%================= 43
\newpage
\Savol{43}
Muntazam 12-burçak tomoni $a=2R\sin\dfrac{180^\circ}{12}=2R\sin15^\circ$, \ $\sin15^\circ=\dfrac{\sqrt6-\sqrt2}4$.

a) $\sqrt6-\sqrt2=2R\cdot\dfrac{\sqrt6-\sqrt2}4\Rightarrow R=2$.

b) Markazdan 12 ta teng yonli uçburçak, uçidagi burçagi $30^\circ$:
\[S=12\cdot\frac12R^2\sin30^\circ=12\cdot\frac12\cdot4\cdot\frac12=12.\]
\Javob{a) 2; \quad b) 12}

%================= 44
\newpage
\Savol{44}
\fig{\begin{tikzpicture}[scale=0.7]
\coordinate (A) at (0,0); \coordinate (B) at (4,0); \coordinate (C) at (5.6,1.6); \coordinate (D) at (1.6,1.6);
\coordinate (O) at (2.8,0.8); \coordinate (S) at (2.8,3.8); \coordinate (K) at (2.8,0);
\draw[thick] (A)--(B)--(C); \draw[thick,dashed] (C)--(D)--(A);
\draw[thick] (S)--(A); \draw[thick] (S)--(B); \draw[thick] (S)--(C); \draw[thick,dashed] (S)--(D);
\draw[ax] (S)--(O); \draw[axs] (S)--(K); \draw[ax] (O)--(K);
\fill (O) circle (1.4pt) node[right]{\scriptsize$O$};
\node[below left] at (A){\small$A$}; \node[below] at (B){\small$B$}; \node[right] at (C){\small$C$}; \node[above left] at (D){\small$D$}; \node[above] at (S){\small$S$}; \node[below] at (K){\scriptsize$K$};
\node[aux] at (3.2,0.45){\tiny$\sqrt3$};
\end{tikzpicture}}
Yon yoqlar asosga bir xil burçak ostida $\Rightarrow$ balandlik asosi $O$ — içki çizilgan doira markazi, $OK=r=\sqrt3$.\\
Romb balandligi $2r=2\sqrt3$, tomoni $a=\dfrac{2\sqrt3}{\sin30^\circ}=4\sqrt3$, \ $S_{\text{asos}}=a\cdot2r=4\sqrt3\cdot2\sqrt3=24$.

a) Yon yoqlar proyeksiyasi — asos: $S_{\text{yon}}=\dfrac{S_{\text{asos}}}{\cos60^\circ}=48$, \ $S_{\text{töla}}=24+48=72$.

b) $SO=OK\tg60^\circ=\sqrt3\cdot\sqrt3=3$, \ $V=\dfrac13\cdot24\cdot3=24$.
\Javob{a) 72; \quad b) 24}

%================= 45
\newpage
\Savol{45}
Işçilarning unumdorligi (1 soatda bajaradigan uluşi) $p_1,\dots,p_5$; \ $s=p_1+p_3$:
\[s+p_2=\frac2{15},\qquad s+p_5=\frac15,\qquad s+p_4=\frac16,\qquad p_2+p_4+p_5=\frac14.\]
Dastlabki uçtasini qöşib, törtinçisini ayiramiz: \ $3s=\dfrac2{15}+\dfrac15+\dfrac16-\dfrac14=\dfrac{8+12+10-15}{60}=\dfrac14\Rightarrow s=\dfrac1{12}$.

a) $p_2=\dfrac2{15}-\dfrac1{12}=\dfrac{8-5}{60}=\dfrac1{20}$ $\Rightarrow$ 20 soat.

b) $p_1+\dots+p_5=s+(p_2+p_4+p_5)=\dfrac1{12}+\dfrac14=\dfrac13$ $\Rightarrow$ 3 soat.
\Javob{a) 20; \quad b) 3}

\ShowJavoblar
\end{document}
